\documentclass{report}
\usepackage{amsmath}
\usepackage{amssymb}
\usepackage{hyperref}
\usepackage[margin=1in]{geometry}
\setlength{\parindent}{0pt}

\begin{document}
\thispagestyle{empty}
\begin{center} \LARGE Assignment 6 \end{center}

\section*{Part 1: Theory and By-Hand Computation}

\subsection*{Q1. Bellman equations on a small MDP}

(a)

\begin{enumerate}
    \item \begin{equation}
        Q^*(A,left)=1.0(5+0.5V^*(end))=5
    \end{equation}
    
    \begin{equation}
        Q^*(A,right)=0.8(2+0.5V^*(B))+0.2(2+0.5V^*(A))
    \end{equation}
    
    \begin{equation}
        V^*(A)=max{Q^*(A,left),Q^*(A,right)}
    \end{equation}
    \item 
    \begin{equation}
        Q^*(B,left)=1+0.5V^*(A)
    \end{equation}
    \begin{equation}
        Q^*(B,right)=0.6(3+0.5V^*(C))+0.4(10)
    \end{equation}
    \begin{equation}
        V^*(B)=max{Q^*(B,left),Q^*(B,right)}
    \end{equation}
    \item 
    \begin{equation}
        Q^*(C,left)=1+0.5V^*(B)
    \end{equation}  

    \begin{equation}
        Q^*(C,right)=8
    \end{equation}

    \begin{equation}
        V^*(C)=max{Q^*(C,left),Q^*(C,right)}
    \end{equation}

\end{enumerate}





(b)
\begin{enumerate}
    \item \begin{equation}
        V^*(C)=max(1+0.5V^*(B),8)
    \end{equation}
    lets just assume 8 is correct. move on with calc
    \item \begin{equation}
        Q^*(B,left)=1+0.5V^*(A)
    \end{equation}
    \begin{equation}
        Q^*(B,right)=8.2
    \end{equation}
    assume 
    \begin{equation}
        V^*(B)=8.2
    \end{equation}
    \item 
    \begin{equation}
        Q^*(A,left)=5
    \end{equation}
    \begin{equation}
        Q^*(A,right)=0.8(2+0.5(8.2))+0.2(2+0.5V^*(A))
    \end{equation}
    if right is optimal,
    \begin{equation}
        V^*(A)=5.28+0.1V^*(A)
    \end{equation}
\end{enumerate}

we can check and see that right is in fact optimal

checking for C, the other option is worse than 8 so 8 was in fact optimal.

and we get

\begin{equation}\begin{split}
        V^*(A) &\approx 5.867\\
        V^*(B) &= 8.2\\
        V^*(C) &= 8\\
\end{split}
\end{equation}

(c)

optimal policy for all states is right


(d)

if \(\gamma\) is 1, then A can return to itself with some probability when going right, and thus \(V^*(A)\) goes to infinity.

\subsection*{Q2. Value iteration by hand on a grid world}

(a)
\begin{equation}
    Q((3,3),E)=0.8(1)+0.1(-0.04)+0.1(-0.04) = -.792
\end{equation}

\begin{equation}
    Q((3,2),N)=0.8(-0.04)+0.1(-1)+0.1(-0.04) = -.136
\end{equation}
\begin{equation}
    Q((3,1),E)=-0.04
\end{equation}

(b)

\begin{equation}
    Q((3,3),E)=0.8(1+0)+0.1(-0.04+0.792)+0.1(-0.04-0.04) = 0.8672
\end{equation}

\begin{equation}
    Q((3,2),W)=0.8(-0.04-0.04)+0.1(-0.04+0.792)+0.1(-0.04-0.04) = 0.0032
\end{equation}

state \((3,1)\) just takes another penalty and the result is \(-.08\)

(c)


\begin{equation}
    \begin{matrix}
        \rightarrow & \rightarrow & \rightarrow &+1\\
        \uparrow & WALL & \uparrow & -1\\
        \uparrow & \rightarrow & \uparrow & \leftarrow
    \end{matrix}
\end{equation}

the nodes in column 3 prefer to finish faster than to avoid the -1.

\subsection*{Q3. The effect of discount and living reward}

(a)

\begin{enumerate}
    \item a small \(\gamma\) means the agent may avoid the tile next to -1 since it is punished. 
    \item with a large \(\gamma\) the agent may decide to more aggressively head towards the +1 tile.
\end{enumerate}

(b)

\begin{enumerate}
    \item with a small \(r\), the agent will try to avoid -1 since existence is punished less
    \item with a large \(r\), the agent will not mind passing by -1 since -1 once is not as bad as -2 each extra step. 
\end{enumerate}

(c)
\(\gamma<1\) is required for cycles to not diverge in reward. the infinite recursive sum might go to \(+\infty\) if the discount factor is 1. 


\section*{Reflection + Statistics Section}
(a)

(b) 

(c) 

(d) 

Contributions: \\



Total: \\

Part 1: \\
Q1: \\

Part 2: \\
Q1: \\

Part 3: \\
Q1: \\

(e)


\end{document}
